% !TITLE 浪淘沙·帘外雨潺潺
% !AUTHOR 李煜
% !DYNASTY 五代十国
% !GENRE 词
% !ORDER 5

\begin{poem}
帘外雨潺潺\textsuperscript{1}，春意阑珊\textsuperscript{2}。罗衾\textsuperscript{3}不耐五更\textsuperscript{4}寒。
梦里不知身是客\textsuperscript{5}，一晌\textsuperscript{6}贪欢\textsuperscript{7}。

独自莫凭栏\textsuperscript{8}，无限江山。别时容易见时难。
流水落花春去也，天上人间\textsuperscript{9}。
\end{poem}

\begin{annotations}
\notetext{1}{潺潺：流水声，这里形容雨声连绵不断。雨声滴答不停，增添了孤独感。}
\notetext{2}{阑珊：衰残、将尽。春意阑珊，春天快要过去了。春将尽、雨连绵——气氛凄凉。}
\notetext{3}{罗衾：丝绸被子。罗，一种轻软的丝织品。即便盖着丝被，也抵不住五更天的寒冷。}
\notetext{4}{五更：天快亮的时候，约凌晨四点到五点。一夜中最冷的时刻，也是人最脆弱的时候。}
\notetext{5}{身是客：自己是异乡之客。李煜亡国后被俘到汴京（开封），实际上是被囚禁的俘虏。说"客"已是美化。}
\notetext{6}{一晌：片刻、一会儿。梦里还不知自己已是阶下囚，贪恋着那片刻的欢愉。}
\notetext{7}{贪欢：贪图欢娱。"一晌贪欢"四个字，是李煜对自己亡国前生活的沉痛忏悔——曾经以为可以永远快乐，原来只是一晌。}
\notetext{8}{凭栏：靠着栏杆（远望）。不要独自靠在栏杆上望远方——因为望见的江山已经不再属于自己了。}
\notetext{9}{天上人间：流水带走了落花，春天一去不返——就像从天堂跌落到人间。也有解释为：曾经的欢乐在天上，如今的痛苦在人间。}
\end{annotations}
